\ifdefined\gkver\else\def\gkver{student}\fi
\documentclass[\gkver]{gaokaozhenti}
\begin{document}

\chapter{2010年北京理综}
%% paperType: 理综物理部分

\begin{enumerate}
\item
%% number: 13
%% paperName: 2010年北京理综第 13 题 6 分
%% typeId: 1
%% type: 单选题
%% chapter: 相对论简介
%% point: 相对论
%% method: 无
%% score: 6
%% degree: 850
%% duplicateId: 0
%% body:
属于狭义相对论基本假设的是：在不同参考系中，\xzanswer{A}
\fourchoices[answer=A]
{真空中光速不变}
{时间间隔具有相对性}
{物体的质量不变}
{物体的能量与质量成正比}
%% answer: A
%% memo:
\memoanswer{
狭义相对论两个基本假设：相对性原理、光速不变原理。光速不变原理：真空中的光速 $c$ 是对任何惯性参照系都适用的普适常量。A 正确。
}

\item
%% number: 14
%% paperName: 2010年北京理综第 14 题 6 分
%% typeId: 1
%% type: 单选题
%% chapter: 光学
%% point: 光的电磁说
%% method: 无
%% score: 6
%% degree: 650
%% duplicateId: 0
%% body:
对于红、黄、绿、蓝四种单色光，下列表述正确的是\xzanswer{C}
\fourchoices[answer=C]
{在相同介质中，绿光的折射率最大}
{红光的频率最高}
{在相同介质中，蓝光的波长最短}
{黄光光子的能量最小}
%% answer: C
%% memo:
\memoanswer{
红、黄、绿、蓝四种单色光的频率依次增大，光从真空进入介质频率不变，B 错。由色散现象同一介质对频率大的光有大的折射率，A 错。频率大的光在真空中和介质中的波长都小，蓝光的波长最短，C 正确。频率大，光子能量大，D 错。
}

\item
%% number: 15
%% paperName: 2010年北京理综第 15 题 6 分
%% typeId: 1
%% type: 单选题
%% chapter: 原子和原子核
%% point: 质能方程
%% method: 无
%% score: 6
%% degree: 650
%% duplicateId: 0
%% body:
太阳因核聚变释放出巨大的能量，同时其质量不断减少。太阳每秒钟辐射出的能量约为 $4\times10^{26}\UJ$，根据爱因斯坦质能方程，太阳每秒钟减少的质量最接近\xzanswer{D}
\fourchoices[answer=D]
{$10^{36}\Ukg$}
{$10^{18}\Ukg$}
{$10^{13}\Ukg$}
{$10^{9}\Ukg$}
%% answer: D
%% memo:
\memoanswer{
根据爱因斯坦的质能方程，$\Delta m=\frac{\Delta E}{c^{2}}=\frac{4\times10^{26}\UJ}{(3\times10^{8}\Ums)^{2}}=4.4\times10^{9}\Ukg$，选项 D 正确。
}

\item
%% number: 16
%% paperName: 2010年北京理综第 16 题 6 分
%% typeId: 1
%% type: 单选题
%% chapter: 曲线运动
%% point: 人造地球卫星
%% method: 无
%% score: 6
%% degree: 450
%% duplicateId: 0
%% body:
一物体静置在平均密度为 $\rho$ 的球形天体表面的赤道上。已知万有引力常量为 $G$，若由于天体自转使物体对天体表面压力恰好为零，则天体自转周期为\xzanswer{D}
\fourchoices[answer=D]
{$\sqrt{\frac{4\pi}{3G\rho}}$}
{$\sqrt{\frac{3}{4\pi G\rho}}$}
{$\sqrt{\frac{\pi}{G\rho}}$}
{$\sqrt{\frac{3\pi}{G\rho}}$}
%% answer: D
%% memo:
\memoanswer{
球形天体表面的赤道上，物体对天体表面压力恰好为零，说明天体对物体的万有引力恰好等于物体随天体运动所需的向心力，有 $G\frac{\rho\frac{4}{3}\pi R^{3}m}{R^{2}}=m\left(\frac{2\pi}{T}\right)^{2}R$，解得 $T=\sqrt{\frac{3\pi}{\rho G}}$。正确选项为 D。
}

\item
%% number: 17
%% paperName: 2010年北京理综第 17 题 6 分
%% typeId: 1
%% type: 单选题
%% chapter: 机械振动
%% point: 振动图象
%% method: 无
%% score: 6
%% degree: 450
%% duplicateId: 0
%% body:
一列横波沿 $x$ 轴正向传播，$a$、$b$、$c$、$d$ 为介质中沿波传播方向上四个质点的平衡位置。某时刻的波形如图 \ref{2010北京理综17a} 所示，此后，若经过 $3/4$ 周期开始计时，则图 \ref{2010北京理综17b} 描述的是\xzanswer{B}
\twopicture[w=6cm, numA=1, numB=2, labelA=2010北京理综17a, labelB=2010北京理综17b]
{figs/17a.png}
{figs/17b.png}
\fourchoices[answer=B]
{$a$ 处质点的振动图象}
{$b$ 处质点的振动图象}
{$c$ 处质点的振动图象}
{$d$ 处质点的振动图象}
%% answer: B
%% memo:
\memoanswer{
由波的图像经过 $3/4$ 周期 $a$ 到达波谷，$b$ 达到平衡位置向下运动，$c$ 到达波峰，$d$ 达到平衡位置向上运动，这四个质点在 $3/4$ 周期开始计时时刻的状态只有 $b$ 符合振动图像。选项 B 正确。
}

\item
%% number: 18
%% paperName: 2010年北京理综第 18 题 6 分
%% typeId: 1
%% type: 单选题
%% chapter: 电路
%% point: 电容
%% method: 无
%% score: 6
%% degree: 650
%% duplicateId: 0
%% body:
用控制变量法，可以研究影响平行板电容器电容的因素（如图）。设两极板正对面积为 $S$，极板间的距离为 $d$，静电计指针偏角为 $\theta$。实验中，极板所带电荷量不变，若\xzanswer{A}
\onepicture[align=right, w=5.5cm]{figs/18.png}
\fourchoices[answer=A]
{保持 $S$ 不变，增大 $d$，则 $\theta$ 变大}
{保持 $S$ 不变，增大 $d$，则 $\theta$ 变小}
{保持 $d$ 不变，减小 $S$，则 $\theta$ 变小}
{保持 $d$ 不变，减小 $S$，则 $\theta$ 不变}
%% answer: A
%% memo:
\memoanswer{
由平行板电容器 $C=\frac{\varepsilon S}{4\pi kd}$ 及 $C=\frac{Q}{U}$，保持 $S$ 不变，增大 $d$，电容 $C$ 减小，电荷量 $Q$ 不变，电势差 $U$ 增大，静电计指针偏角 $\theta$ 增大。保持 $d$ 不变，减小 $S$，电容 $C$ 减小，电荷量 $Q$ 不变，电势差 $U$ 增大，静电计指针偏角 $\theta$ 增大。正确选项 A。
}

\item
%% number: 19
%% paperName: 2010年北京理综第 19 题 6 分
%% typeId: 1
%% type: 单选题
%% chapter: 电磁感应
%% point: 自感与互感，涡流
%% method: 无
%% score: 6
%% degree: 450
%% duplicateId: 0
%% body:
在如图所示的电路中，两个相同的小灯泡 $\mathrm{L}_{1}$ 和 $\mathrm{L}_{2}$ 分别串联一个带铁芯的电感线圈 $L$ 和一个滑动变阻器 $R$。闭合开关 $S$ 后，调整 $R$，使 $\mathrm{L}_{1}$ 和 $\mathrm{L}_{2}$ 发光的亮度一样，此时流过两个灯泡的电流均为 $I$。然后，断开 $S$。若 $t'$ 时刻再闭合 $S$，则在 $t'$ 前后的一小段时间内，正确反映流过 $\mathrm{L}_{1}$ 的电流 $i_{1}$、流过 $\mathrm{L}_{2}$ 的电流 $i_{2}$ 随时间 $t$ 变化的图像是\xzanswer{B}
\onepicture[align=right, w=5.5cm]{figs/19.png}
\fourchoices[answer=B, ispicture=true, h=2.2cm]
{figs/19a.png}
{figs/19b.png}
{figs/19c.png}
{figs/19d.png}
%% answer: B
%% memo:
\memoanswer{
由电路实物图可得，与滑动变阻器 $R$ 串联的 $\mathrm{L}_{2}$，没有自感直接变亮，电流 $i_{2}$ 变化图像如 A 中图线，C、D 错误。带铁芯的电感线圈串联的 $\mathrm{L}_{1}$，由于自感强电流逐渐变大，B 正确。
}

\item
%% number: 20
%% paperName: 2010年北京理综第 20 题 6 分
%% typeId: 1
%% type: 单选题
%% chapter: 运动和力
%% point: 动量定理
%% method: 无
%% score: 6
%% degree: 450
%% duplicateId: 0
%% body:
如图，若 $x$ 轴表示时间，$y$ 轴表示位置，则该图像反映了某质点做匀速直线运动时，位置与时间的关系。若令 $x$ 轴和 $y$ 轴分别表示其它的物理量，则该图像又可以反映在某种情况下，相应的物理量之间的关系。下列说法中正确的是\xzanswer{C}
\onepicture[align=right, w=4cm]{figs/20.png}
\fourchoices[answer=C]
{若 $x$ 轴表示时间，$y$ 轴表示动能，则该图像可以反映某物体受恒定合外力作用做直线运动过程中，物体动能与时间的关系}
{若 $x$ 轴表示频率，$y$ 轴表示动能，则该图像可以反映光电效应中，光电子最大初动能与入射光频率之间的关系}
{若 $x$ 轴表示时间，$y$ 轴表示动量，则该图像可以反映某物在沿运动方向的恒定合外力作用下，物体动量与时间的关系}
{若 $x$ 轴表示时间，$y$ 轴表示感应电动势，则该图像可以反映静置于磁场中的某闭合回路，当磁感应强度随时间均匀增大时，增长合回路的感应电动势与时间的关系}
%% answer: C
%% memo:
\memoanswer{
根据动量定理 $P-P_{0}=Ft$ 得 $P=P_{0}+Ft$，说明动量和时间是线性关系，纵截距为初动量，C 正确。结合 $P=\sqrt{2mE_{k}}$ 得 $\sqrt{2mE_{k}}=Ft+P_{0}$，说明动能和时间是抛物线，A 错误。根据光电效应方程 $E_{km}=h\nu-W$，说明最大初动能和入射光频率是线性关系，但截距为负值，B 错误。当磁感应强度随时间均匀增大时，增长合回路的磁通量均匀增大，根据法拉第电磁感应定律增长合回路的感应电动势等于磁通量的变化率，是一个定值不随时间变化，D 错误。
}

\item
%% number: 21
%% paperName: 2010年北京理综第 21 题 18 分
%% typeId: 4
%% type: 实验题
%% chapter: 电路
%% point: 电路实验
%% method: 无
%% score: 18
%% degree: 450
%% duplicateId: 0
%% body:
（1）甲同学要把一个量程为 $200\UuA$ 的直流电流计 G，改装成量程是 $0\sim4\UV$ 的直流电压表。

\begin{steps}
	\item
	按图 \ref{2010北京理综21a} 所示电路、用半偏法测定电流计 G 的内电阻 $r_{g}$，其中电阻 $R_{0}$ 约为 $1\UkO$。为使 $r_{g}$ 的测量值尽量准确，在以下器材中，电源 $E$ 应选用\tkanswer[1cm]{B}，电阻器 $R_{1}$ 应选用\tkanswer[1cm]{C}，电阻器 $R_{2}$ 应选用\tkanswer[1cm]{F}（选填器材前的字母）。

	（A）电源（电动势 $1.5\UV$）　　　　（B）电源（电动势 $6\UV$）

	（C）电阻箱（$0\sim999.9\UO$）　　　（D）滑动变阻器（$0\sim500\UO$）

	（E）电位器（一种可变电阻，与滑动变阻器相当）（$0\sim5.1\UkO$）

	（F）电位器（$0\sim51\UkO$）

	\item
	该同学在开关断开情况下，检查电路连接无误后，将 $R_{2}$ 的阻值调至最大。后续的实验操作步骤依次是：\tkanswer[1cm]{B}，\tkanswer[1cm]{C}，\tkanswer[1cm]{A}，\tkanswer[1cm]{E}，最后记录 $R_{1}$ 的阻值并整理好器材。（请按合理的实验顺序，选填下列步骤前的字母）

	（A）闭合 $S_{1}$

	（B）闭合 $S_{2}$

	（C）调节 $R_{2}$ 的阻值，使电流计指针偏转到满刻度

	（D）调节 $R_{2}$ 的阻值，使电流计指针偏转到满刻度的一半

	（E）调节 $R_{1}$ 的阻值，使电流计指针偏转到满刻度的一半

	（F）调节 $R_{1}$ 的阻值，使电流计指针偏转到满刻度

	\item
	如果所得的 $R_{1}$ 的阻值为 $300.0\UO$，则图 \ref{2010北京理综21a} 中被测电流计 G 的内阻 $r_{g}$ 的测量值为\tkanswer[1.5cm]{300.0} $\UO$，该测量值\tkanswer[1.3cm]{略小于}实际值（选填“略大于”、“略小于”或“等于”）。

	\item
	给电流计 G\tkanswer[0.8cm]{串}联（选填“串”或“并”）一个阻值为\tkanswer[1.3cm]{19.7} $\UkO$ 的电阻，就可以将该电流计 G 改装为量程 $4\UV$ 的电压表。
\end{steps}

\onepicture[align=right, w=6cm, num=1, label=2010北京理综21a]{figs/21a.png}

（2）乙同学要将另一个电流计 G 改装成直流电压表，但他仅借到一块标准电压表 $V_{0}$、一个电池组 $E$、一个滑动变阻器 $R'$ 和几个待用的阻值准确的定值电阻。

①该同学从上述具体条件出发，先将待改装的表 G 直接与一个定值电阻 $R$ 相连接，组成一个电压表；然后用标准电压表 $V_{0}$ 校准。请你画完图 \ref{2010北京理综21b} 中方框中的校准电路图。

\onepicture[align=right, h=4.2cm, num=2, label=2010北京理综21b]{figs/21b.png}

②实验中，当定值电阻 $R$ 选用 $17.0\UkO$ 时，调整滑动变阻器 $R'$ 的阻值，电压表 $V_{0}$ 的示数是 $4.0\UV$ 时，表 G 的指针恰好指到满量程的五分之二；当 $R$ 选用 $7.0\UkO$ 时，调整 $R'$ 的阻值，电压表 $V_{0}$ 的示数是 $2.0\UV$，表 G 的指针又指到满量程的五分之二。

由此可以判定，表 G 的内阻 $r_{g}$ 是\tkanswer[1.2cm]{3.0} $\UkO$，满偏电流 $I_{g}$ 是\tkanswer[1.2cm]{0.50} $\UmA$。若要将表 G 改装为量程是 $15\UV$ 的电压表，应配备一个\tkanswer[1.2cm]{27} $\UkO$ 的电阻。

%% answer: 
\jdanswer{
（1）①B、C、F；②B、C、A、E；③$300.0$，略小于；④串，$19.7$。

（2）①如图所示；②$3.0$，$0.50$，$27$。
}
%% memo:
\memoanswer{
（1）①半偏法测量表头内阻时，首先选择滑动变阻器（必须大于电路所需的最小电阻），根据电路的电压为电动势，电路的最大电流为表头的满偏电流，则最小电阻为 $\frac{6\UV}{0.2\times10^{-3}\UA}=3.0\times10^{4}\UO=30\UkO$，或 $\frac{1.5\UV}{0.2\times10^{-3}\UA}=7.5\times10^{4}\UO=75\UkO$，考虑到保护电阻 $1\UkO$，则可知调节滑动变阻器使表头满偏时滑动变阻器的阻值分别接近 $29\UkO$ 或 $6.5\UkO$，电路图中 $R_{2}$ 是滑动变阻器，不能选 D 和 E，只能选 F。表头满偏时滑动变阻器的阻值越大，实验的误差越小。所以电源选择电动势为 $6\UV$ 的 B，而且滑动变阻器 F 的阻值也满足调节所需。而 $R_{1}$ 是用来测量表头内阻的电阻箱，只能选 C。

②实验步骤：第一步闭合 $S_{2}$，选 B；第二步调节 $R_{2}$ 阻值使电流计满偏，选 C；第三步闭合 $S_{1}$，选 A；第四步调节 $R_{1}$ 阻值使电流计半偏，选 E；第五步读出 $R_{1}$ 阻值为待测表头的内阻。

③$R_{1}$ 的示数为待测表头的内阻是 $300.0\UO$，闭合 $S_{1}$ 后，电路的总电阻减小，当表头半偏时干路上的电流就大于表头的满偏电流，流过电阻箱的电流就大于表头的满偏电流，所以电阻箱的阻值略小于表头的内阻。

④给表头串联一个电阻可以改装为电压表，改装后的电压表的内阻为 $R_{V}=\frac{U_{m}}{I_{g}}=\frac{4\UV}{0.2\times10^{-3}\UA}=2.0\times10^{4}\UO=20\UkO$，则串联电阻的大小为 $20\UkO-300\UO=19.7\UkO$。

（2）①本实验是对改装电压表进行校准的，将改装电压表和标准电压表并联后要求电压从 0 至满偏变化，所以滑动变阻器采用分压接法，电路连接如图。

\onepicture[h=4.2cm]{figs/21b_an.png}

②表头的刻度均匀，当指针恰好指到满量程的五分之二时，流过电流表的电流为 $0.4I_{g}$，根据欧姆定律分别有 $4\UV=0.4I_{g}(r_{g}+17\times10^{3}\UO)$ 和 $2\UV=0.4I_{g}(r_{g}+7\times10^{3}\UO)$，解得 $r_{g}=3.0\UkO$，$I_{g}=0.50\UmA$。

量程为 $15\UV$ 的改装电压表的内阻 $R_{V}=\frac{U_{m}}{I_{g}}=\frac{15\UV}{0.5\times10^{-3}\UA}=3.0\times10^{4}\UO=30\UkO$，所以改装时串联的电阻是 $30\UkO-3.0\UkO=27\UkO$。
}

\item
%% number: 22
%% paperName: 2010年北京理综第 22 题 16 分
%% typeId: 6
%% type: 计算题
%% chapter: 功和能
%% point: 机械能守恒定律
%% method: 无
%% score: 16
%% degree: 650
%% duplicateId: 0
%% body:
如图，跳台滑雪运动员经过一段加速滑行后从 O 点水平飞出，经过 $3.0\Us$ 落到斜坡上的 A 点。已知 O 点是斜坡的起点，斜坡与水平面的夹角 $\theta=37^{\circ}$，运动员的质量 $m=50\Ukg$。不计空气阻力。（取 $\sin37^{\circ}=0.60$，$\cos37^{\circ}=0.80$；$g$ 取 $10\Umsq$）求：
\begin{enumerate}
	\item
	A 点与 O 点的距离 $L$；
	\item
	运动员离开 O 点时的速度大小；
	\item
	运动员落到 A 点时的动能。
\end{enumerate}
\onepicture[align=right, w=6cm]{figs/22.png}
%% answer: 
\jdanswer{
\begin{enumerate}
	\item
	$75\Um$
	\item
	$20\Ums$
	\item
	$32500\UJ$
\end{enumerate}
}
%% memo:
\memoanswer{
（1）运动员在竖直方向做自由落体运动，有 $L\sin37^{\circ}=\frac{1}{2}gt^{2}$，A 点与 O 点的距离 $L=\frac{gt^{2}}{2\sin37^{\circ}}=75\Um$。

（2）设运动员离开 O 点的速度为 $v_{0}$，运动员在水平方向做匀速直线运动，即 $L\cos37^{\circ}=v_{0}t$，解得 $v_{0}=\frac{L\cos37^{\circ}}{t}=20\Ums$。

（3）由机械能守恒，取 A 点为重力势能零点，运动员落到 A 点时的动能为 $E_{kA}=mgh+\frac{1}{2}mv_{0}^{2}=32500\UJ$。
}

\item
%% number: 23
%% paperName: 2010年北京理综第 23 题 18 分
%% typeId: 6
%% type: 计算题
%% chapter: 磁场
%% point: 带电粒子在磁场中的运动
%% method: 无
%% score: 18
%% degree: 450
%% duplicateId: 0
%% body:
利用霍尔效应制作的霍尔元件以及传感器，广泛应用于测量和自动控制等领域。

如图 \ref{2010北京理综23a}，将一金属或半导体薄片垂直置于磁场 $B$ 中，在薄片的两个侧面 a、b 间通以电流 $I$ 时，另外两侧 c、f 间产生电势差，这一现象称为霍尔效应。其原因是薄片中的移动电荷受洛伦兹力的作用向一侧偏转和积累，于是 c、f 间建立起电场 $E_{H}$，同时产生霍尔电势差 $U_{H}$。当电荷所受的电场力与洛伦兹力处处相等时，$E_{H}$ 和 $U_{H}$ 达到稳定值，$U_{H}$ 的大小与 $I$ 和 $B$ 以及霍尔元件厚度 $d$ 之间满足关系式 $U_{H}=R_{H}\frac{IB}{d}$，其中比例系数 $R_{H}$ 称为霍尔系数，仅与材料性质有关。

\threepicture[w=4.5cm, numA=1, numB=2, numC=3, labelA=2010北京理综23a, labelB=2010北京理综23b, labelC=2010北京理综23c]
{figs/23a.png}
{figs/23b.png}
{figs/23c.png}

\begin{enumerate}
	\item
	设半导体薄片的宽度（c、f 间距）为 $l$，请写出 $U_{H}$ 和 $E_{H}$ 的关系式；若半导体材料是电子导电的，请判断图 \ref{2010北京理综23a} 中 c、f 哪端的电势高；
	\item
	已知半导体薄片内单位体积中导电的电子数为 $n$，电子的电荷量为 $e$，请导出霍尔系数 $R_{H}$ 的表达式。（通过横截面积 $S$ 的电流 $I=nevS$，其中 $v$ 是导电电子定向移动的平均速率）；
	\item
	图 \ref{2010北京理综23b} 是霍尔测速仪的示意图，将非磁性圆盘固定在转轴上，圆盘的周边等距离地嵌装着 $m$ 个永磁体，相邻永磁体的极性相反。霍尔元件置于被测圆盘的边缘附近。当圆盘匀速转动时，霍尔元件输出的电压脉冲信号图像如图 \ref{2010北京理综23c} 所示。

	（A）若在时间 $t$ 内，霍尔元件输出的脉冲数目为 $P$，请导出圆盘转速 $N$ 的表达式。

	（B）利用霍尔测速仪可以测量汽车行驶的里程。除此之外，请你展开“智慧的翅膀”，提出另一个实例或设想。
\end{enumerate}
%% answer: 
\jdanswer{
\begin{enumerate}
	\item
	$U_{H}=E_{H}l$，c 端电势高。
	\item
	$\frac{1}{ne}$。
	\item
	$N=\frac{P}{mt}$，提出的实例或设想合理即可。
\end{enumerate}
}
%% memo:
\memoanswer{
（1）$U_{H}=E_{H}l$，c 端电势高。

（2）由 $U_{H}=R_{H}\frac{IB}{d}$①，得 $R_{H}=U_{H}\frac{d}{IB}=E_{H}l\frac{d}{IB}$②。当电场力与洛伦兹力相等时 $eE_{H}=evB$，得 $E_{H}=vB$③。又 $I=nevS$④。将③④代入②得 $R_{H}=vBl\frac{d}{IB}=vl\frac{d}{nevS}=\frac{ld}{neS}=\frac{1}{ne}$。

（3）（A）由于在时间 $t$ 内，霍尔元件输出的脉冲数目为 $P$，则 $P=mNt$，圆盘转速为 $N=\frac{P}{mt}$。

（B）提出的实例或设想合理即可。
}

\item
%% number: 24
%% paperName: 2010年北京理综第 24 题 20 分
%% typeId: 6
%% type: 计算题
%% chapter: 运动和力
%% point: 动量和能量
%% method: 无
%% score: 20
%% degree: 250
%% duplicateId: 0
%% body:
雨滴在穿过云层的过程中，不断与漂浮在云层中的小水珠相遇并结合为一体，其质量逐渐增大。现将上述过程简化为沿竖直方向的一系列碰撞。已知雨滴的初始质量为 $m_{0}$，初速度为 $v_{0}$，下降距离 $l$ 后与静止的小水珠碰撞且合并，质量变为 $m_{1}$。此后每经过同样的距离 $l$ 后，雨滴均与静止的小水珠碰撞且合并，质量依次变为 $m_{2}$、$m_{3}$……$m_{n}$……（设各质量为已知量）。不计空气阻力。
\begin{enumerate}
	\item
	若不计重力，求第 $n$ 次碰撞后雨滴的速度 $v_{n}'$；
	\item
	若考虑重力的影响，

	（A）求第 1 次碰撞前、后雨滴的速度 $v_{1}$ 和 $v_{1}'$；

	（B）求第 $n$ 次碰撞后雨滴的动能 $\frac{1}{2}m_{n}v_{n}'^{2}$。
\end{enumerate}
%% answer: 
\jdanswer{
\begin{enumerate}
	\item
	$\frac{m_{0}}{m_{n}}v_{0}$
	\item
	$v_{1}=\sqrt{v_{0}^{2}+2gl}$，$v_{1}'=\frac{m_{0}}{m_{1}}\sqrt{v_{0}^{2}+2gl}$；$\frac{1}{2m_{n}}\left[m_{0}^{2}v_{0}^{2}+2gl(m_{0}^{2}+m_{1}^{2}+\cdots+m_{n}^{2})\right]$
\end{enumerate}
}
%% memo:
\memoanswer{
（1）不计重力，全过程中动量守恒，$m_{0}v_{0}=m_{n}v_{n}'$，得 $v_{n}'=\frac{m_{0}}{m_{n}}v_{0}$。

（2）若考虑重力的影响，雨滴下降过程中做加速度为 $g$ 的匀加速运动，碰撞瞬间动量守恒。

（A）第 1 次碰撞前 $v_{1}^{2}=v_{0}^{2}+2gl$，$v_{1}=\sqrt{v_{0}^{2}+2gl}$。第 1 次碰撞后 $m_{0}v_{1}=m_{1}v_{1}'$，$v_{1}'=\frac{m_{0}}{m_{1}}v_{1}=\frac{m_{0}}{m_{1}}\sqrt{v_{0}^{2}+2gl}$①。

（B）第 2 次碰撞前 $v_{2}^{2}=v_{1}'^{2}+2gl$，利用①化简得 $v_{2}^{2}=\left(\frac{m_{0}}{m_{1}}\right)^{2}v_{0}^{2}+\frac{m_{0}^{2}+m_{1}^{2}}{m_{1}^{2}}\cdot2gl$②。第 2 次碰撞后利用②得 $v_{2}'^{2}=\left(\frac{m_{1}}{m_{2}}\right)^{2}v_{2}^{2}=\left(\frac{m_{0}}{m_{2}}\right)^{2}v_{0}^{2}+\frac{m_{0}^{2}+m_{1}^{2}}{m_{2}^{2}}\cdot2gl$。同理，第 3 次碰撞后 $v_{3}'^{2}=\left(\frac{m_{0}}{m_{3}}\right)^{2}v_{0}^{2}+\frac{m_{0}^{2}+m_{1}^{2}+m_{2}^{2}}{m_{3}^{2}}\cdot2gl$，……第 $n$ 次碰撞后 $v_{n}'^{2}=\left(\frac{m_{0}}{m_{n}}\right)^{2}v_{0}^{2}+\frac{\sum_{i=0}^{n-1}m_{i}^{2}}{m_{n}^{2}}\cdot2gl$，动能 $\frac{1}{2}m_{n}v_{n}'^{2}=\frac{1}{2m_{n}}\left(m_{0}^{2}v_{0}^{2}+2gl\sum_{i=0}^{n-1}m_{i}^{2}\right)$。
}

\end{enumerate}

\end{document}
